\chapter{Options~: parité et grecques}\label{ch:iv2:options-parity-and-greeks}

\begin{ivtrains}
Ce que rapporte une position d'options, ce qu'aucun prix ne peut violer, ce que mesure chaque grecque et où elle est la plus grande, et à quoi ressemble la journée d'un hedger. Sauf indication contraire de l'énoncé, les options sont européennes sur une action ne versant pas de dividende, un contrat porte sur 100 actions et les grecques sont celles du modèle de Black--Scholes.
\end{ivtrains}

\section{Profils de gain et bornes}

\begin{question}[$\star$ \iqroles{QT, Trader, QR, Risque} \iqfirm{toutes} \iqlevel{junior}]\label{iq:iv2:options-parity-and-greeks:1}
Vous achetez un call de prix d'exercice 100 pour 5~\$ et vendez un call de prix d'exercice 110 pour 2~\$, sur la même action et avec la même échéance. Quels sont votre gain maximal, votre perte maximale et votre point mort à l'échéance~? Quel scénario de marché cela traduit-il~?
\end{question}

\begin{question}[$\star$ \iqroles{QT, Trader, QR, Risque} \iqfirm{toutes} \iqlevel{junior}]\label{iq:iv2:options-parity-and-greeks:2}
Vous achetez un call et un put, tous deux de prix d'exercice 100, pour 8~\$ au total. Où se situent vos points morts à l'échéance~? Que gagnez-vous ou perdez-vous si l'action termine à 105, puis à 120~? Sur quoi pariez-vous~?
\end{question}

\begin{question}[$\star$ \iqroles{QT, Trader, QR} \iqfirm{teneur de marché, hedge fund, banque} \iqlevel{junior}]\label{iq:iv2:options-parity-and-greeks:3}
Vous achetez un call 90, vendez deux calls 100 et achetez un call 110, en payant 2~\$ au total. Quels sont votre profit maximal, votre perte maximale et où se situent vos points morts à l'échéance~? Pourquoi cette position ne peut-elle jamais avoir un prix négatif~?
\end{question}

\begin{question}[$\star$ \iqroles{QT, Trader, QR} \iqfirm{teneur de marché, hedge fund, banque} \iqlevel{junior}]\label{iq:iv2:options-parity-and-greeks:4}
Les taux sont nuls. Vous achetez un call et vendez un put de même prix d'exercice 100, pour un coût net de 2~\$. Quelle position cela constitue-t-il~? Quel doit être le cours de l'action pour qu'il n'y ait pas d'arbitrage~? Quel est votre profit si l'action termine à 110~?
\end{question}

\begin{question}[$\star$ \iqroles{QT, Trader, QR} \iqfirm{teneur de marché, hedge fund, banque} \iqlevel{junior}]\label{iq:iv2:options-parity-and-greeks:5}
Une action cote 100, le prix d'exercice est 90, le temps jusqu'à l'échéance est d'un an et le taux est de 5~\%, capitalisé en continu. Quel est le prix minimal possible du call européen~? Il se négocie à 12~\$. Donnez les transactions qui permettent de prendre l'arbitrage, et le gain qu'elles dégagent.
\end{question}

\begin{question}[$\star$ \iqroles{QT, Trader, QR} \iqfirm{teneur de marché, hedge fund, banque} \iqlevel{junior}]\label{iq:iv2:options-parity-and-greeks:6}
Le taux est de 5~\%, capitalisé en continu. Un put européen d'un an, de prix d'exercice 100, se négocie à 97~\$. Pourquoi est-il trop cher, quoi que fasse l'action~? Quel est le prix maximal qu'il peut atteindre~? Quelles transactions permettent de verrouiller le profit~?
\end{question}

\begin{question}[$\star$ \iqroles{QT, Trader, QR} \iqfirm{teneur de marché, hedge fund} \iqlevel{junior}]\label{iq:iv2:options-parity-and-greeks:7}
Sur une même action et une même échéance, un put de prix d'exercice 100 coûte 6~\$, et un put de prix d'exercice 110 coûte 5~\$. Les taux sont nuls. Qu'est-ce qui cloche~? Que négociez-vous, et quel est votre pire scénario~? Quel est l'écart maximal que pourraient présenter les deux puts~?
\end{question}

\begin{question}[$\star$ \iqroles{QT, Trader, QR, Risque} \iqfirm{toutes} \iqlevel{junior}]\label{iq:iv2:options-parity-and-greeks:8}
Pour un call acheté et un put acheté, quels sont les signes du delta, du gamma, du vega et du theta dans le modèle de Black--Scholes à taux nuls~? Lesquels sont identiques pour les deux, et pourquoi doivent-ils l'être~?
\end{question}

\begin{question}[$\star$ \iqroles{QT, Trader, Risque} \iqfirm{toutes} \iqlevel{junior}]\label{iq:iv2:options-parity-and-greeks:9}
Vous êtes acheteur de 10 contrats de call de delta 0.6, vendeur de 5 contrats de put de delta $-0.3$, et vendeur de 200 actions. Chaque contrat porte sur 100 actions. Quel est votre delta total, en actions~? Combien d'actions devez-vous négocier pour le rendre nul~?
\end{question}

\begin{question}[$\star$ \iqroles{QT, Trader} \iqfirm{teneur de marché, hedge fund, banque} \iqlevel{junior}]\label{iq:iv2:options-parity-and-greeks:10}
Vous êtes acheteur de 5 contrats de call à la monnaie, chacun portant sur 100 actions, de delta 0.5 et de gamma 0.02 par dollar. Combien d'actions devez-vous vendre pour vous couvrir~? L'action monte de 1~\$. Quel est le nouveau delta, et que faites-vous~?
\end{question}

\begin{question}[$\star$ \iqroles{QT, Trader, Risque} \iqfirm{teneur de marché, hedge fund, banque} \iqlevel{junior}]\label{iq:iv2:options-parity-and-greeks:11}
Votre book affiche un gamma de 1\,000 actions par dollar et un theta de $-\$1\,500$ par jour, et il est couvert en delta. Grosso modo, que gagnez-vous ou perdez-vous si l'action bouge de 2~\$ aujourd'hui~? Et pour 1~\$~? Quel mouvement quotidien faut-il pour atteindre le point mort~?
\end{question}

\begin{question}[$\star$ \iqroles{QT, Trader, Risque} \iqfirm{teneur de marché, hedge fund, banque} \iqlevel{junior}]\label{iq:iv2:options-parity-and-greeks:12}
Votre book a un vega de 500~\$ par point de volatilité. La volatilité implicite de l'échéance la plus proche monte de 3 points, et celle des échéances éloignées de 1 point~; le vega du book est de 300~\$ sur les échéances proches et de 200~\$ sur les éloignées. Quel est votre profit~? Que gagneriez-vous si vous utilisiez simplement le vega total et un mouvement unique de 3 points~?
\end{question}

\begin{question}[$\star$ \iqroles{QT, QR, Trader} \iqfirm{teneur de marché, hedge fund, banque} \iqlevel{junior}]\label{iq:iv2:options-parity-and-greeks:13}
Un call et un put à la monnaie d'un an portent sur une action valant 100~\$, avec une volatilité de 20~\% et des taux nuls. De combien chacun varie-t-il, grosso modo, si le taux monte d'un point de pourcentage~? Lequel des deux prend de la valeur, et pourquoi~?
\end{question}

\begin{question}[$\star$ \iqroles{QT, QR, Trader} \iqfirm{teneur de marché, hedge fund, banque} \iqlevel{junior}]\label{iq:iv2:options-parity-and-greeks:14}
Pourquoi ne devez-vous jamais exercer par anticipation un call américain sur une action ne versant aucun dividende, lorsque les taux sont positifs~? Démontrez-le à l'aide d'une borne. Pourquoi le même raisonnement échoue-t-il pour un put~?
\end{question}

\section{La parité sous questions}

\begin{question}[$\star\star$ \iqroles{QT, Trader, QR} \iqfirm{teneur de marché, hedge fund, banque} \iqlevel{junior}]\label{iq:iv2:options-parity-and-greeks:15}
Une action cote 100, les options d'un an ont un prix d'exercice de 100, et le taux est de 5~\%, capitalisé en continu, si bien que la valeur actuelle du prix d'exercice est d'environ 95.12. Le call cote 6~\$ et le put 3~\$. Quelle devrait être la différence~? Donnez les transactions qui rapportent de l'argent, leur coût aujourd'hui, et leur résultat pour chaque cours de l'action à l'échéance.
\end{question}

\begin{question}[$\star\star$ \iqroles{QT, Trader, QR} \iqfirm{teneur de marché, hedge fund} \iqlevel{junior, confirmé}]\label{iq:iv2:options-parity-and-greeks:16}
Un call moins un put d'un an, de prix d'exercice 100, se négocie à 3~\$, sur une action cotant 100 qui ne verse pas de dividende, la valeur actuelle du prix d'exercice étant de 95.12. La parité donne 4.88~\$. Vous essayez de verrouiller l'écart en achetant le call, en vendant le put et en vendant l'action à découvert, et votre courtier affiche des frais d'emprunt de 2~\% par an. Quel effet ont ces frais~? Quels frais d'emprunt l'écart implique-t-il~? Pourquoi cet écart vous renseigne-t-il généralement sur l'emprunt plutôt que sur un profit facile~?
\end{question}

\begin{question}[$\star\star$ \iqroles{QT, QR} \iqfirm{teneur de marché, hedge fund, banque} \iqlevel{junior, confirmé}]\label{iq:iv2:options-parity-and-greeks:17}
Pour des options américaines sur une action ne versant aucun dividende, la parité se transforme en un couple d'inégalités. Avec une action à 100, un strike de 100, un taux de 5~\% capitalisé en continu et un an jusqu'à l'échéance, entre quelles bornes doit se situer le call moins le put~? Démontrez la borne supérieure.
\end{question}

\begin{question}[$\star\star$ \iqroles{QT, Trader, QR} \iqfirm{teneur de marché, hedge fund, banque} \iqlevel{junior, confirmé}]\label{iq:iv2:options-parity-and-greeks:18}
Vous achetez un call et vendez un put de prix d'exercice 90, puis vendez un call et achetez un put de prix d'exercice 110, tous à échéance d'un an. Que vaut cette position à l'échéance, quel que soit le cours final de l'action~? Le marché demande 18.70~\$ pour elle. Quel taux d'intérêt, capitalisé en continu, cela implique-t-il~? Si vous pouvez emprunter à 5~\%, que faites-vous~?
\end{question}

\begin{question}[$\star\star$ \iqroles{QT, QR} \iqfirm{teneur de marché, hedge fund} \iqlevel{junior}]\label{iq:iv2:options-parity-and-greeks:19}
Sur une action ne versant pas de dividende et avec des taux positifs, un call européen d'un mois de prix d'exercice 100 est coté 4~\$, et un call de trois mois est coté 3.50~\$. S'agit-il d'un arbitrage~? Qu'en est-il si un dividende important intervient entre-temps~? Quelle est la transaction à réaliser si c'en est un~?
\end{question}

\begin{question}[$\star\star$ \iqroles{QT, QR, Risque} \iqfirm{teneur de marché, hedge fund, banque} \iqlevel{junior, confirmé}]\label{iq:iv2:options-parity-and-greeks:20}
Une action de 100~\$ a une volatilité de 20~\%, et les taux sont nuls. Quel est le gamma d'un call à la monnaie avec un an à courir, puis avec un seul jour de bourse restant, sachant qu'il y en a 252 par an~? Combien faut-il s'éloigner du strike, en dollars, pour réduire de moitié le gamma avec un jour restant~? Pourquoi une option vendue à court terme effraie-t-elle les gens~?
\end{question}

\begin{question}[$\star\star$ \iqroles{QT, QR} \iqfirm{teneur de marché, hedge fund, banque} \iqlevel{junior, confirmé}]\label{iq:iv2:options-parity-and-greeks:21}
Un put européen de prix d'exercice 100 a un an à courir sur une action de 50~\$ avec une volatilité de 20~\% et un taux de 5~\%, capitalisé en continu. Quelle est sa valeur, et gagne-t-il ou perd-il de la valeur si un jour de plus s'écoule, toutes choses égales par ailleurs~? Combien cela fait-il par an~? Quand le theta peut-il être positif~?
\end{question}

\begin{question}[$\star\star$ \iqroles{QT, Trader, QR} \iqfirm{teneur de marché, hedge fund} \iqlevel{junior}]\label{iq:iv2:options-parity-and-greeks:22}
Une action cote 100, avec une volatilité de 20~\%, des taux nuls et un an jusqu'à l'échéance. Quel est le vega par point de volatilité pour des strikes de 80, 100 et 120~? Pourquoi n'est-il pas symétrique autour de 100~?
\end{question}

\begin{question}[$\star\star$ \iqroles{QT, QR, Trader} \iqfirm{teneur de marché, hedge fund} \iqlevel{junior}]\label{iq:iv2:options-parity-and-greeks:23}
Une action de 100~\$ a une volatilité de 40~\%, des taux nuls et une option d'un an de prix d'exercice 100. Quels sont le delta de l'option et la probabilité risque-neutre qu'elle termine dans la monnaie~? Un trader affirme que le delta est la probabilité. De combien cette affirmation se trompe-t-elle ici~? Dans quels cas est-elle juste~?
\end{question}

\begin{question}[$\star\star$ \iqroles{QT, QR} \iqfirm{teneur de marché, hedge fund, banque} \iqlevel{junior}]\label{iq:iv2:options-parity-and-greeks:24}
Un call d'un an sur une action de rendement en dividende de 2~\% a un delta de 0.55. Quel est le delta du put de même strike~? Dérivez-le. Quel serait-il sans rendement en dividende~?
\end{question}

\begin{question}[$\star\star$ \iqroles{QT, Trader, QR} \iqfirm{teneur de marché, hedge fund} \iqlevel{junior}]\label{iq:iv2:options-parity-and-greeks:25}
Vous détenez un straddle de gamma constant égal à 0.04 par dollar, couvert en delta à 100. L'action monte à 103 puis revient à 100, et vous rééquilibrez votre couverture à chacun de ces points. Quelles sont vos transactions, et quel est votre gain par straddle~? Vérifiez-le à l'aide de la formule du gamma. Quel doit être le theta d'une journée pour que cela atteigne le point mort~?
\end{question}

\begin{question}[$\star\star$ \iqroles{QT, Trader, Risque} \iqfirm{teneur de marché, hedge fund, banque} \iqlevel{junior, confirmé}]\label{iq:iv2:options-parity-and-greeks:26}
Ce matin, votre livre a un delta de 200 actions, un gamma de 50 actions par dollar, un vega de 300~\$ par point et un theta de $-\$150$ par jour. Aujourd'hui, l'action monte de 2~\$ et la volatilité implicite baisse d'un point. Que auriez-vous dû gagner, ligne par ligne~? Vous avez en réalité gagné 120~\$~: qu'est-ce qui reste inexpliqué, et que pourrait-il l'expliquer~?
\end{question}

\begin{question}[$\star\star$ \iqroles{QT, Trader, QR} \iqfirm{teneur de marché, hedge fund} \iqlevel{junior}]\label{iq:iv2:options-parity-and-greeks:27}
Vous vendez un put 90 à 3.00~\$ et achetez un call 110 à 2.50~\$. Leurs deltas valent $-0.27$ et $0.35$, et leurs vegas 0.31 et 0.37 par point. Que gagnez-vous à l'échéance si le sous-jacent vaut 80, 100 et 120~? Quels sont le delta et le vega de votre position~? Que gagnez-vous si la volatilité implicite du call 110 monte d'un point et celle du put 90 baisse d'un point~?
\end{question}

\begin{question}[$\star\star$ \iqroles{QT, Trader, Analyste} \iqfirm{toutes} \iqlevel{junior}]\label{iq:iv2:options-parity-and-greeks:28}
Les taux sont nuls. Vous détenez une action à 100, achetez un put 90 à 3~\$ et vendez un call 110 à 3~\$. Quels sont votre gain et votre perte maximaux~? En utilisant la parité put-call, montrez que la position est identique à un call spread augmenté d'un montant de liquidités. Combien doit valoir le call spread 90-110~?
\end{question}

\begin{question}[$\star\star$ \iqroles{QT, Trader, Analyste} \iqfirm{toutes} \iqlevel{junior}]\label{iq:iv2:options-parity-and-greeks:29}
Les taux sont nuls. Vous achetez une action à 100 et vendez un call 105 à 3~\$. Quel est votre profit à 85, 105 et 120~? Quelle position plus simple a le même profil de gains, et à quel prix doit-elle se négocier~? Pourquoi a-t-on tort de penser que vendre des calls contre une action rapporte un «~revenu~» sans risque~?
\end{question}

\begin{question}[$\star\star$ \iqroles{QT, QR, Trader} \iqfirm{banque, teneur de marché, hedge fund} \iqlevel{junior}]\label{iq:iv2:options-parity-and-greeks:30}
Un contrat à terme vaut 100. Des options européennes sur ce contrat ont un prix d'exercice de 95, expirent dans six mois, et le taux est de 5~\%, capitalisé en continu. Quelle est la différence entre le call et le put~? Démontrez-le. Pourquoi la formule diffère-t-elle de celle des options sur action~?
\end{question}

\begin{question}[$\star\star$ \iqroles{QT, QR, Trader} \iqfirm{teneur de marché, hedge fund, banque} \iqlevel{junior, confirmé}]\label{iq:iv2:options-parity-and-greeks:31}
Une option verse 1~\$ si l'action termine au-dessus de 100, et rien sinon. L'action vaut 100 avec une volatilité de 20~\%, et les taux sont nuls. Quelle est sa valeur dans le modèle de Black--Scholes, et comment l'obtient-on à partir du prix des calls ordinaires~? Quel est le problème du couvreur à l'approche de l'échéance~?
\end{question}

\begin{question}[$\star\star$ \iqroles{QT, QR, Risque} \iqfirm{teneur de marché, hedge fund} \iqlevel{junior}]\label{iq:iv2:options-parity-and-greeks:32}
Vous détenez un call spread 100-110. Où son gamma est-il positif, où est-il négatif, et approximativement où s'annule-t-il~? Pourquoi~? Que cela implique-t-il pour votre couverture lorsque l'action monte dans l'intervalle~?
\end{question}

\begin{question}[$\star\star$ \iqroles{QT, QR, Trader} \iqfirm{teneur de marché, hedge fund} \iqlevel{junior, confirmé}]\label{iq:iv2:options-parity-and-greeks:33}
Une action à 100~\$ a des options à un an de prix d'exercice 100, avec des taux nuls. La volatilité implicite du call est de 25~\%, celle du put de 27~\%. En dollars, quel est l'écart de parité entre les deux options~? Quel est l'emprunt implicite~? Pourquoi la plupart des systèmes n'affichent-ils qu'une volatilité implicite par strike~?
\end{question}

\begin{question}[$\star\star$ \iqroles{QT, QR} \iqfirm{teneur de marché, hedge fund, banque} \iqlevel{junior, confirmé}]\label{iq:iv2:options-parity-and-greeks:34}
Un put américain de prix d'exercice 100 porte sur une action valant 1~\$, avec un an à courir et un taux de 5~\%, capitalisé en continu. Faut-il exercer maintenant~? Montrez-le à l'aide d'une borne sur la valeur que peut atteindre le put européen. Quel rôle joue un taux positif~? Avec un taux nul, exerceriez-vous jamais~?
\end{question}

\section{Couverture et cas plus difficiles}

\begin{question}[$\star\star\star$ \iqroles{QT, QR, Trader} \iqfirm{teneur de marché, hedge fund, banque} \iqlevel{junior, confirmé}]\label{iq:iv2:options-parity-and-greeks:35}
Vous avez vendu une option digitale de 1 million de dollars~: elle verse 1 million de dollars si l'action est au-dessus de 100 à l'échéance, et rien sinon. L'action vaut 100 avec une volatilité de 20~\%, et les taux sont nuls. Quelle est votre couverture en delta, en actions, avec un an devant vous, puis avec un seul jour de bourse restant~? Un dealer évalue en général ce produit comme un call spread~: comment, et qu'apporte-t-il à la couverture~?
\end{question}

\begin{question}[$\star\star\star$ \iqroles{QT, Trader, Risque} \iqfirm{teneur de marché, hedge fund, banque} \iqlevel{junior, confirmé}]\label{iq:iv2:options-parity-and-greeks:36}
C'est le dernier jour. Vous êtes vendeur de 1\,000 calls de prix d'exercice 100, portant chacun sur 100 actions, et acheteur de 100\,000 actions en couverture. L'action clôture à 100.02~\$, et les détenteurs peuvent exercer après la clôture, en totalité ou pas du tout, et vous ne le saurez que le lendemain. Pendant la soirée, l'action chute à 98~\$. Que gagnez-vous ou perdez-vous dans chacun des deux cas~? Comment appelle-t-on ce risque~? Qu'auriez-vous dû faire~?
\end{question}

\begin{question}[$\star\star\star$ \iqroles{QT, QR, Trader} \iqfirm{teneur de marché, hedge fund, banque} \iqlevel{junior, confirmé}]\label{iq:iv2:options-parity-and-greeks:37}
Un call américain de prix d'exercice 100 porte sur une action valant 110~\$ qui verse 2~\$ demain. Après cela, aucun autre dividende n'est attendu, et l'option expire dans un trimestre. Le taux est de 5~\%, capitalisé en continu. Un put européen de même prix d'exercice et de même échéance, après le dividende, vaut 0.40~\$. Faut-il exercer aujourd'hui~? Démontrez une règle qui compare le dividende à ce que l'on abandonne.
\end{question}

\begin{question}[$\star\star\star$ \iqroles{QT, QR} \iqfirm{teneur de marché, hedge fund, banque} \iqlevel{junior, confirmé}]\label{iq:iv2:options-parity-and-greeks:38}
Une action à 100 monte de 20~\% ou baisse de 20~\% à chacune de deux périodes. Le taux par période est de 5~\%, simple. Évaluez un put à deux périodes de prix d'exercice 100, européen puis américain. Où exercez-vous par anticipation, et quelle est la prime de ce droit~?
\end{question}

\begin{question}[$\star\star\star$ \iqroles{QT, QR, Trader} \iqfirm{teneur de marché, hedge fund} \iqlevel{junior, confirmé}]\label{iq:iv2:options-parity-and-greeks:39}
Vous détenez un call à un mois de prix d'exercice 100, couvert chaque jour, avec une volatilité implicite de 30~\% et des taux nuls. Son gamma vaut 0.046 par dollar lorsque l'action est à 100, et 0.0038 lorsqu'elle est à 120. Au bout d'un mois, l'action fluctue de 2~\% par jour dans les deux cas, mais dans l'un elle reste proche de 100, et dans l'autre elle bondit tôt vers 120 et y reste. Approximativement, à combien s'élève le gain de gamma d'une journée type dans chacun des cas~? Que cela nous apprend-il sur la volatilité réalisée comme mesure de votre profit~?
\end{question}

\begin{question}[$\star\star\star$ \iqroles{QT, QR, Trader} \iqfirm{teneur de marché, hedge fund} \iqlevel{junior, confirmé}]\label{iq:iv2:options-parity-and-greeks:40}
Une action vaut 100, avec une volatilité de 20~\%, des taux nuls et un an jusqu'à l'échéance. Vous vendez un straddle at-the-money et achetez assez de strangles 90-110, un put 90 plus un call 110, pour avoir un vega nul à 20~\%. Combien de strangles~? Quelle prime encaissez-vous~? Quel est votre profit si toutes les volatilités implicites passent simultanément à 25~\% ou à 15~\%, sans rien d'autre qui change~? Comment s'appelle la position longue que vous détenez~?
\end{question}

\begin{question}[$\star\star\star$ \iqroles{QR, QT} \iqfirm{teneur de marché, hedge fund} \iqlevel{junior, confirmé}]\label{iq:iv2:options-parity-and-greeks:41}
Vous vendez un call à un an at-the-money sur une action valant 100~\$ avec une volatilité de 20~\% et des taux nuls, et vous le couvrez 252 fois, avec une volatilité égale à la volatilité implicite. Quel est l'écart-type de votre résultat~? Donnez l'approximation et vérifiez-la par simulation. De combien diminue-t-il si vous couvrez quatre fois plus souvent~? Qu'est-ce que cela implique pour la tarification~?
\end{question}

\begin{question}[$\star\star\star$ \iqroles{QR, QT} \iqfirm{teneur de marché, hedge fund} \iqlevel{confirmé}]\label{iq:iv2:options-parity-and-greeks:42}
Vous couvrez un call à un an at-the-money vendu sur une action valant 100~\$, avec une volatilité de 20~\%, des taux nuls, à chacun des 252 jours de bourse, en payant 0.1~\% de la valeur sur chaque action échangée, soit 0.2~\% aller-retour. De combien faut-il relever la volatilité pour couvrir ces coûts, selon la règle de Leland~? Quel est ce coût en dollars de prime~? Vérifiez-le par une simulation des coûts de rééquilibrage.
\end{question}

\begin{question}[$\star\star\star$ \iqroles{QT, QR, Trader} \iqfirm{teneur de marché, hedge fund} \iqlevel{junior, confirmé}]\label{iq:iv2:options-parity-and-greeks:43}
Une action de 100~\$ a des options at-the-money d'un an et d'un mois, avec une volatilité de 20~\% pour les deux, et des taux nuls. Vous êtes vendeur d'un lot à un an. Combien de lots à un mois achetez-vous pour avoir un vega nul~? Quel est alors le gamma net~? Si la volatilité à un mois varie environ $\sqrt{12}$ fois autant que celle à un an, combien de lots à un mois faut-il plutôt~?
\end{question}

\begin{question}[$\star\star\star$ \iqroles{QR, QT} \iqfirm{teneur de marché, hedge fund} \iqlevel{junior, confirmé}]\label{iq:iv2:options-parity-and-greeks:44}
Un digital verse 1~\$ au-dessus de 100 dans un an, sur une action de volatilité 20~\% et des taux nuls. Une hausse de la volatilité le rend-elle plus cher~? Calculez la variation par point pour des cours de 90, 100 et 110, et trouvez le cours auquel le signe s'inverse. Pourquoi ce n'est-il pas 100~? Qu'a un call ordinaire à la place~?
\end{question}

\begin{question}[$\star\star\star$ \iqroles{QT, Trader, Analyste} \iqfirm{teneur de marché, hedge fund} \iqlevel{junior, confirmé}]\label{iq:iv2:options-parity-and-greeks:45}
Les teneurs de marché, pris ensemble, sont estimés être acheteurs de gamma à hauteur de 1 million d'actions par dollar sur une action, à l'approche d'une échéance et près d'un strike important. Le titre échange 5 millions d'actions par jour. S'il monte de 1~\$, que font les teneurs de marché, et en quelle proportion du volume~? Et s'ils sont vendeurs de la même quantité~? Que devient le comportement du cours dans chacun des cas~? Qu'est-ce qui rend cette estimation peu fiable~?
\end{question}

\begin{question}[$\star\star\star$ \iqroles{QR, QT} \iqfirm{teneur de marché, hedge fund} \iqlevel{confirmé}]\label{iq:iv2:options-parity-and-greeks:46}
Supposons que la valeur d'un call $C(S, K)$ soit homogène de degré un en le cours de l'action et le prix d'exercice. Montrez que $C = S\,\partial_S C + K\,\partial_K C$. Utilisez cette relation pour obtenir un delta à partir des seuls prix de marché, avec des taux nuls, lorsque les calls de strikes 99 et 101 sur une action à 100~\$ valent 8.44 et 7.52, et que le call de strike 100 vaut 7.97. Comparez-le au delta de Black--Scholes à 20~\% de volatilité pour un call à un an. Quelle hypothèse sur le smile ce résultat suppose-t-il~?
\end{question}

\begin{question}[$\star\star\star$ \iqroles{QT, QR} \iqfirm{teneur de marché, hedge fund} \iqlevel{junior, confirmé}]\label{iq:iv2:options-parity-and-greeks:47}
Les taux sont nuls, et des calls sur une même action et une même échéance sont cotés~: 90 à 14.00~\$, 100 à 8.00~\$, 110 à 5.50~\% et 120 à 2.00~\$. Quelles conditions ces prix vérifient-ils, et lesquelles violent-ils~? Quelle est votre transaction, et quels sont vos pires et meilleurs résultats à l'échéance~? Que dit le tableau sur la probabilité de terminer au-dessus de 95, au-dessus de 105 et au-dessus de 115~?
\end{question}

\begin{question}[$\star\star\star\star$ \iqroles{QR} \iqfirm{hedge fund, teneur de marché, banque} \iqlevel{confirmé}]\label{iq:iv2:options-parity-and-greeks:48}
Montrez comment répliquer tout gain $f(S_T)$ deux fois dérivable à l'aide d'une obligation, d'un forward et d'un strip de calls et de puts, autour d'un prix de référence $K_0$. Appliquez-le à $f = S_T^2$ avec une action à 100~\$ de volatilité 20~\% dans le modèle de Black--Scholes, des taux nuls et une échéance d'un an, avec $K_0 = 100$. Combien vaut le strip, et combien devrait-il valoir~? Quel est le poids de chaque option~?
\end{question}

\begin{question}[$\star\star\star\star$ \iqroles{QR, QT} \iqfirm{hedge fund, teneur de marché} \iqlevel{confirmé}]\label{iq:iv2:options-parity-and-greeks:49}
Vous vendez un call at-the-money de trois mois sur une action valant 100~\$ avec une volatilité implicite de 30~\%, des taux nuls, et l'action réalise une volatilité de 20~\%. Quel est votre profit espéré~? Comparez une couverture par deltas à la volatilité implicite à une couverture par deltas à la volatilité réalisée. Pour chacune, quelle est la formule du profit le long d'une trajectoire, et quelles sont l'espérance et l'écart-type du profit avec des couvertures quotidiennes~? Lequel choisiriez-vous si vous connaissiez les 20~\%, et pourquoi la plupart des desks utilisent-ils l'implicite~?
\end{question}

\begin{question}[$\star\star\star\star$ \iqroles{QR} \iqfirm{hedge fund, banque} \iqlevel{confirmé}]\label{iq:iv2:options-parity-and-greeks:50}
Un put américain perpétuel de prix d'exercice 100 porte sur une action de volatilité 20~\%, avec un taux de 5~\% capitalisé en continu et sans dividende. Déterminez sa valeur et le cours de l'action auquel on exerce, en résolvant l'équation qu'il doit satisfaire. Montrez que la règle d'exercice est celle qui maximise la valeur. Combien vaut-il avec l'action à 100~? Quelle est la limite si le taux tend vers zéro~?
\end{question}
